\documentclass[11pt]{article}
\usepackage[margin=1in]{geometry}
\usepackage{amsmath,amssymb,hyperref}
\title{Forecast visibility, errors, and reversible preparation}
\date{}
\begin{document}
\maketitle

\paragraph{Result.}
A warning that reaches the crew authorized to install a shield can enable
preparation but need not help the separate crew authorized to divert.
At the central cost point of the \href{https://github.com/eternal-november/eos-research-programmes/pull/355}{original maintenance model}, early warning to both roles gives worst mean cost $13/32$;
early warning confined to the preparation gate gives $7/16$ even with
a standing installation permit. In a separate matched extension,
the ability to remove an installed shield after an advance confirmation
reduces cost from $3/8$ to $21/64$. In the latter comparison the initial
installation cost is still sunk and the capacity it used is not returned
retroactively.

\section{Physical law and case scope}
Keep the pinned prior model, including two already accepted loads, gate
authority, the continuous relay, an optional shielding installation using
$C_2$, and the independent detour switch $E_3$. The old obligations are
scheduled from actual eligibility and intake evidence through the
\href{https://github.com/eternal-november/eos-research-programmes/pull/353}{exact source scheduler}; they are not spare
capacities asserted by the forecast. Write $F=1$ when the existing
matching law permits installation. If intake damage puts $C_2$ and
$D_2$ in competition with mandatory jobs, installation is unavailable.
If the detour instead requires $C_3$, it conflicts with the inherited
tick-3 relay and closure. Every feasible intervention in this study
preserves timely discharge of the old loads, coverage at ticks 0--3,
exclusive slots, and new delivery by tick 7.

The weather potential $H\in\{0,1\}$ disrupts the main zone at tick 5,
not the physically distinct detour. Without shielding a disruption
produces a safe one-tick hold; with shielding it produces none.
No control rule assigns a realized main-zone delay to a detour run.
All changed policies and all costs are calculated from interventions
on the same exogenous tapes, with the original joint laws A and B and
two actual old-health completions on each tape. The baseline warning
$W$ is available at its source before the installation cut (tick 2).
The mission crew receives trusted clearance by tick 2. Either the
mission crew or a gate acting on a standing permit decides installation
at that phase. The two are never given simultaneous installation
authority. The mission crew decides the route at tick 3. If the gate
controls installation, the mission crew learns whether a shield was
installed at tick 3, not the private warning bit itself. Received
messages precede their named decision phase. The result is conditional
on exact clearance, authenticated receipt and the fixed timing
convention; no real warning-data or hardware performance is inferred.

\section{Comparable policies, lead and recipient}
Charge one unit for the reactive hold, $c$ for installation, and $d$ for
diversion. In this comparison a deployed shield is retained without
additional charge. One strategy stays and uses the safe reactive hold
if needed. A second always installs wherever lawfully clear and otherwise
stays; an instruction to install when physically barred is not a fair
policy. A third installs on a positive warning received by the
authorized decision maker by tick 2 and otherwise stays. A fourth
diverts when a positive warning reaches the mission crew by the switch
phase, without installing. Always divert where feasible is a further
control. Every policy is evaluated on the same route and resource law
and warning availability; different strategy classes deliberately do not
use some of that common information. The unrestricted comparison can
combine installation and avoidance contingent on observations.

Under the base early-mission information pattern and at $c=3/8,d=3/4$,
the worst means over A and B are
\[
\begin{array}{l|r}
\text{Reactive stay}&1/2\\
\text{Always install when clear}&7/16\\
\text{Install only on a warning}&15/32\\
\text{Divert only on a warning}&1/2\\
\text{Always divert}&3/4\\
\text{Best lawful combination}&13/32
\end{array}
\]
A timely warning has value for coordinating these choices; it does not
make a warning-only rule preferable to unconditional lawful installation
in this cost regime. Across both laws the direct minimax optimization
reproduces the prior best value. At the same costs, delivery to the mission
crew at tick 1 or before the decision in tick 2 gives $13/32$. Delivery
at tick 3 (only diversion can still respond), delivery at tick 4 (neither
can), no delivery, or a detour blocked by inherited relay work gives
$7/16$ in each case. The equal late values are specific to the central cost
point; the two timings do not have identical control surfaces.
Giving the gate the warning at tick 1 permits its standing installation
rule to condition on it, but when the mission crew receives no warning,
it cannot make the same warning-directed switch on blocked tapes.
The best gate-only value is $7/16$, with or without the standing permit,
and giving the warning to both with standing authority attains $13/32$.
This comparison includes the installed-status observation on each tape;
the gate's knowledge is not treated as automatically shared.

\section{Asymmetric warning errors}
Put a separate draw on every tape. The delivered bit is $V=W$ except
that the channel changes $0$ to $1$ with probability $\alpha$, and
$1$ to $0$ with probability $\beta$. These are conditional on the
forecast $W$, not on the physical hazard. Under each original law
the actual rates against $H$ are
\[
   \Pr(V=0\mid H=1)=(1-\alpha+3\beta)/4,\qquad
   \Pr(V=1\mid H=0)=(1+3\alpha-\beta)/4.
\]
Keep these two rates separate rather than summarizing them by a
nominal accuracy. Under early mission receipt at the center costs,
$(\alpha,\beta)=(0,0)$ gives baseline rates $(1/4,1/4)$ and the
value $13/32$; $(0,1/4)$ gives $(7/16,3/16)$ and value $55/128$;
$(1/4,0)$ gives $(3/16,7/16)$ and value $7/16$. Equal error rates
$(1/2,1/2)$ leave no predictive information and give $7/16$.
At the lower installation cost $1/8$, always installing when clear
costs $5/16$ and is minimax in all nine tested error pairs. This extension
does not assert that arbitrary false-warning processes are independent
of the old-work state; it explicitly assumes independent channel flips
conditional on $W$ and retains the two different original joint laws.

\section{Reversibility after a confirmation}
Now all compared policies share a truthful, separately supplied advance
confirmation of the impending main-zone condition at tick 3, before
either route is traversed. This is an added sensing assumption, not a
forecast delivered sooner or an observed outcome substituted across
interventions. All compared policies use the same original costs and a
new surcharge $g$ for running the mission with the installed shield.
In the irreversible arrangement, after using $C_2$ the mission must keep
the shield and use the main route. In the reversible arrangement,
a release through the route selector at tick 3 costs $r$ and makes the
unshielded main route or the independently available detour admissible.
The shield is physically no longer present on the selected route.
The selector can release without allocating an additional worker slot;
if release requires a resource or conflicts with a relay, this result
does not apply. The installation charge and the completed use of $C_2$
survive release. Thus the comparison holds costs fixed, including any
release fee for an admissible release, and changes the transition
relation, not just the value of $c$.

At $c=1/8,g=1/4,d=3/4$, an already installed shield is kept if the
main hazard is confirmed: retaining is less costly than either exposed
route. If no hazard is confirmed, it is released precisely when
$r<g$ (a tie at equality). Let $u=\min(r,g)$, and for the irreversible
controller put $u=g$. Let $x,y\in[0,1]$ be the installation probabilities
on clear tapes for $W=1,0$ respectively. With no preparation, the
confirmation permits optimal diversion on $H=1$, at mean $d/2$.
The clear tapes under law A have mass $1/4$ each for $(H,W)=(1,1),(0,0)$;
under law B all four pairs have mass $1/8$. This gives
\[
 L_A=\frac38+\frac{-3x/8+(1/8+u)y}{4},\qquad
 L_B=\frac38+\frac{(u-1/4)(x+y)}8 .
\]
For $u\le1/4$, $L_B$ is minimized by $x=y=1$, and this also makes
$L_A=L_B$. The exact minimax value is $5/16+u/4$. At $r=1/16$ it
is $21/64$ versus $3/8$ when retention is mandatory. Reversal with a fee
at least $1/4$ preserves the option but gives no economic benefit on
these tapes. Logged matched runs include retaining, removing and taking
the physically independent detour after removal; the last is feasible
and does not use a disturbance observation as a detour outcome, even
when it is not the minimizing choice. The absence of a benefit at high
fee does not make the two control systems physically equivalent.

\section{Exact checks and limits}
For each law the replay propagates rational tape weights and generates
decisions from actual received data at the two cuts. It enumerates all
admissible installation rules and all continuation choices on reached
observation cells; if confirmation makes each continuation's cost
known identically under both laws, pointwise dominated continuations
can be dropped. The original primal two-policy minimax calculation is
checked independently against the lower envelope of risk under mixture
weights on the two laws, enumerating endpoints and intersections.
The parameter register includes nine timing/recipient/capacity cases,
eighteen asymmetric-error and installation-cost cases, and six reversal
fees checked against the formula above. It exercises 69 distinct
feasible plant route histories and retains actual closure and relay
assertions in every history. All processing is exact and finite;
the closed-form reversal result is proved on its whole stated interval.
The matched direct table proof needs the same information, resources,
and trace laws and gives the same answer. The original study compares
\href{https://arxiv.org/abs/2103.05526}{information-age predictive control}
and distributionally robust optimization; this extension is specifically
about prior commitments, recipient-specific warning utility and the
sunk-cost boundary for withdrawal. It does not establish novelty in
predictive control or a physical test. Together with the original
workload it addresses the written C7 comparisons in this simulated
envelope; a materially different operational plant must discharge these
assumptions afresh.
\end{document}
